\begin{table}[htbp]
\centering
\footnotesize
\caption{Estimator performance in both declared Monte Carlo designs}
\begin{tabularx}{\linewidth}{lp{0.26\linewidth}rrrrX}
\toprule
Design & Estimator & Bias & RMSE & MC SE & Coverage & Wilson 95\% \\
\midrule
Nonlinear & OLS & 0.633 & 0.639 & 0.003 & 0.00 & [0.0000, 0.0038] \\
Nonlinear & In-sample orthogonal ML & 0.007 & 0.050 & 0.002 & 0.94 & [0.9269, 0.9557] \\
Nonlinear & Linear Lasso plug-in & 0.592 & 0.599 & 0.003 & 0.00 & [0.0000, 0.0038] \\
Nonlinear & Cross-fitted DML & 0.039 & 0.072 & 0.002 & 0.92 & [0.9037, 0.9371] \\
Sparse & OLS & 0.001 & 0.051 & 0.002 & 0.94 & [0.9213, 0.9513] \\
Sparse & Linear Lasso plug-in & 0.610 & 0.610 & 0.000 & 0.00 & [0.0000, 0.0038] \\
Sparse & Cross-fitted DML & 0.003 & 0.040 & 0.001 & 0.94 & [0.9213, 0.9513] \\
\bottomrule
\end{tabularx}
\par\vspace{3pt}\begin{minipage}{\linewidth}\footnotesize Source: author's calculations from the documented inputs.\allowbreak{} Coverage refers to nominal 95\% treatment-effect intervals; MC SE is the standard error of mean bias across repetitions.\allowbreak{} Wilson intervals quantify uncertainty in estimated coverage, not treatment effects.\allowbreak{}\end{minipage}
\end{table}
